\documentclass[12pt]{article}
\usepackage{amsmath, amsthm, amssymb}
\usepackage{geometry}
\geometry{margin=1in}
\usepackage{hyperref}

\newtheorem{theorem}{Theorem}
\newtheorem{proposition}[theorem]{Proposition}
\theoremstyle{definition}
\newtheorem{definition}[theorem]{Definition}
\theoremstyle{remark}
\newtheorem{remark}[theorem]{Remark}

\title{Othering:\\
The Grammar of Distance}
\author{Diego Rinc\'on\\
\small Independent}
\date{}

\begin{document}

\maketitle

\begin{abstract}
Every grammar others what it cannot read.
Othering is not hostility --- it is structure:
the formal consequence of two grammars
that share no resonance object \cite{resonance},
no common token, no isomorphism between them.
The Wall \cite{synthesis} is othering crystallized:
the arithmetic others the spectral,
biomass others experience,
the finite others infinity,
and each othering carves a face of the Wall.
Levinas saw it first in ethics:
the face of the Other resists totalization,
cannot be reduced to the Same,
is infinitely other.
The mathematics says the same thing:
the archimedean fiber is infinitely other
to the finite primes.
Resonance \cite{resonance} is the undoing of othering:
the discovery that two grammars were never truly other ---
they were reading the same motive \cite{motives}
from opposite sides.
Self-othering is the hardest Wall:
a system that cannot read itself
has othered its own ground.
\end{abstract}

\section{Othering as Grammar Failure}

\begin{definition}[Othering]
Let $G_1$ and $G_2$ be two grammars \cite{grammar}
--- two systems of reading persistence.
$G_1$ \emph{others} $G_2$ if there is no object $M$
that $G_1$ can read directly from $G_2$'s output:
no resonance \cite{resonance}, no isomorphism,
no shared token that crosses between them.

Formally: $G_1$ others $G_2$
if $\mathrm{Hom}_{G_1}(G_2(-), M) = 0$
for every object $M$ accessible to $G_1$.
The two grammars share no common reading.
They have othered each other into mutual opacity.

Othering is symmetric in structure but not in power:
$G_1$ may other $G_2$ while $G_2$ reads $G_1$
--- the dominated grammar sees the dominant one
but is not seen in return.
The Wall between them is one-sided.
\end{definition}

\begin{proposition}[The Wall is othering]
Every face of the Wall \cite{five} is an othering:
\begin{enumerate}
\item \emph{Arithmetic Wall}: the grammar of primes
  others the grammar of zeros.
  The primes are countable, discrete, arithmetically generated.
  The zeros are spectral, continuous, analytically distributed.
  Neither grammar reads the other directly.
  The explicit formula \cite{week09} holds them in balance
  without either crossing.
\item \emph{Logic Wall (G\"odel)}: the grammar of provability
  others the grammar of truth.
  There are true sentences the provability grammar cannot reach.
  The true sentence has othered itself out of the proof system.
\item \emph{Computation Wall (Turing)}: the grammar of halting
  others the grammar of programs.
  No program grammar can read whether all programs halt.
  The halting predicate is other to every algorithm.
\item \emph{Consciousness Wall}: the grammar of biomass
  others the grammar of experience.
  Neural firing is readable; experience is not.
  The land space \cite{consciousness} is other to the cipher.
\item \emph{Infinity Wall}: the grammar of the finite
  others the grammar of the infinite.
  Every finite grammar hits its ceiling
  and the infinite remains beyond.
\end{enumerate}
Each Wall is an othering that mathematics has not yet undone.
\end{proposition}

\section{Levinas and the Face}

\begin{proposition}[The face resists othering]
Emmanuel Levinas (1961) identified the ethical form of othering:
the reduction of the Other to a category of the Same.
To other a person is to make them legible in your grammar ---
to convert their irreducibility into a role, a type, a function.

The face is what resists this conversion.
The face of the Other is the point at which
your grammar runs out:
you cannot totalize it, categorize it, domesticate it.
It is infinitely other.

The mathematical analogue:
the archimedean place is the face of $\mathrm{Spec}\,\mathbb{Z}$.
The arithmetic grammar (Frobenius at each prime $p$)
runs out at $\infty$.
The face --- the archimedean fiber --- cannot be animated
by the grammar that reads every finite prime.
It remains other.
The Arithmetic Wall is the refusal of $\mathrm{Spec}\,\mathbb{Z}$
to totalize its own face.
\end{proposition}

\section{Self-Othering}

\begin{definition}[Self-othering]
A system \emph{self-others} when it cannot read its own ground:
when the grammar that reads everything else
cannot read the condition of its own reading.

G\"odel's incompleteness is self-othering:
the provability grammar cannot read
its own consistency from inside.
The system has othered its own foundation.

The hard problem of consciousness is self-othering:
the organism reads the world
but cannot read the reading.
Experience is other to the machinery that produces it.
The cipher cannot read its own key.

The Arithmetic Wall is self-othering:
$\mathrm{Spec}\,\mathbb{Z}$ is trying to complete itself
but lacks the fiber at $\infty$ to do it.
The arithmetic grammar cannot read
its own archimedean place.
It has othered its own completion.
\end{definition}

\begin{remark}[The most isolated]
Self-othering is the deepest isolation.
An external Wall separates two systems that might, in principle,
find a resonance token and bridge the gap.
A self-Wall separates a system from itself:
the gap is not between two grammars
but within one grammar, between its surface and its ground.

The hard problem is this:
consciousness is not other to the world
but other to the system that produces it.
The organism cannot step outside its own decryption
to read the decryption.
Only the cipher-holder holds the key,
and the key is the cipher-holder.
The ouroboros of self-othering.
\end{remark}

\section{Resonance as the Undoing of Othering}

\begin{theorem}[Resonance dissolves othering]
\label{thm:undo}
Two grammars $G_1$ and $G_2$ that have othered each other
are un-othered by a resonance object \cite{resonance}:
an object $M$ such that
\[
G_1(M) \cong G_2(M),
\]
--- both grammars read the same thing from $M$.
The discovery of $M$ does not eliminate the Wall
but reveals that it was always transparent:
$G_1$ and $G_2$ were reading the same motive \cite{motives}
from different sides.

Every comparison isomorphism is an un-othering:
$H_{\rm dR}(X) \cong H_{\rm sing}(X)$
says that de Rham cohomology and singular cohomology
had othered each other for a century
and were reading the same motive all along.
The isomorphism does not merge the grammars.
It shows they were never truly other.

The Riemann Hypothesis is the claim
that the zeros and the primes can be un-othered:
that there exists a self-adjoint operator $D$ \cite{week10}
whose spectrum is the zeros
and whose trace formula is the primes,
an object that both grammars read simultaneously.
Proof of RH is proof that arithmetic and spectrum
were never truly other.
\end{theorem}

\begin{remark}[Othering and compassion]
In ethics, compassion is the undoing of othering:
the recognition that the Other's grammar
is not foreign but the same grammar
spoken in a different body.

In mathematics, comparison isomorphisms are compassion:
the formal recognition that two grammars
have been othering the same object
since the beginning.

Both say: we were always reading the same thing.
The Wall was always thinner than it looked.
We did not know how to see through it.
\end{remark}

\begin{thebibliography}{9}

\bibitem{synthesis}
Rinc\'on, D. (2026).
The Arithmetic Wall. Preprint.

\bibitem{resonance}
Rinc\'on, D. (2026).
Resonance: When Two Grammars Read the Same Thing. Preprint.

\bibitem{five}
Rinc\'on, D. (2026).
Five Walls, One Aura. Preprint.

\bibitem{grammar}
Rinc\'on, D. (2026).
Mathematics as Cohomological Grammar. Preprint.

\bibitem{motives}
Rinc\'on, D. (2026).
Motives, Standard Conjectures, and the Unification of $L$-functions.
Preprint.

\bibitem{week09}
Rinc\'on, D. (2026).
Week 09: The Resonance Object. Preprint.

\bibitem{week10}
Rinc\'on, D. (2026).
Week 10: Self-Adjointness. Preprint.

\bibitem{consciousness}
Rinc\'on, D. (2026).
Consciousness: The Land Space. Preprint.

\bibitem{math}
Rinc\'on, D. (2026).
Mathematics. Preprint.

\end{thebibliography}

\end{document}
